\section{Experimental Evaluation}
\label{sec:eval}
The evaluation asks whether one registered-harness path can execute bounded
tasks across two simulator scenes with different signal and protocol mappings
while retaining a checkable task-to-approval-to-receipt-to-sample trace. R20 is
the primary result: six sequential tasks through one registered harness, three
per scene on one simulator instance per scene. It demonstrates this configured
software path, not configuration-reuse benefit or process attainment. R13 is
secondary exploratory protocol and data-path evidence whose exact runtime
source was not pinned. All experiments use repository simulators. An anonymous
evidence archive is supplied as Supplementary File~S1; its trace matrix and
run records support the R20 checks reported below.

\subsection{Evidence Units and Experimental Procedure}
\label{sec:testbed}
\label{sec:procedure}
R20 uses the registered \texttt{omp} harness in cast-film and annealing scenes,
with three sequential task executions per scene on one simulator instance per
scene. R13 uses configured protocol and data paths in the same two simulator
scenes. Earlier experiments are retained separately in Supplementary File~S1
and are not pooled with these results.

\subsection{Registered-Harness Process Path (R20)}
\label{sec:harness-e2e}
R20 ran the local v0.8.0 working-tree application with the registered
\texttt{omp} adapter and GLM-5.3-Flash. Supplementary File~S1 contains the
recorded base revisions, working-tree patches, dependency lockfiles, prompt
assets, and run logs. This is a source-reconstruction kit rather than an
immutable release, and the external model service does not guarantee
deterministic replay. Each of six sequential executions created a separate
line, product, recipe, run, channel, agent, and task context. Three executions
shared one simulator instance per scene. Cast film used OPC UA for screw-speed
control and MQTT for thickness; annealing used Modbus TCP for both the zone-3
setpoint and strip-temperature probe. The model received a live control read,
process measurement, and target band, then selected one value from the
scenario-defined prevalidated set (125/130/135 rpm or 735/740/745~$^{\circ}$C)
for a bounded \texttt{recipe\_trial}; this was a constrained choice, not an
unrestricted optimization search. A synthetic administrator approved each
proposal. Each route log contains one trial and a 1/1 verified device receipt.

After verifying the route ACK and returned trial result, the runner initiated
a 20-s delay, then probed, ingested, and queried one value per task under the
active-run tag through the isolated SQLite time-series fallback. The runner
checked task state after the sample query; all six task records were then
\texttt{COMPLETED}. The runner reported 76 passing checks and no failures, and
the separate trace audit found a complete task-to-sample record in all six
executions: each had a matching approval ID, one readback-verified trial
receipt, and one query result whose run ID and value matched the active run and
explicit probe. The task-completion timestamp was not retained, so no
task-completion-to-sample interval can be computed. Table~\ref{tab:q8} labels
Film-1--3 map to \texttt{cast-film-r1--r3}, and Anneal-1--3 to
\texttt{annealing-r1--r3}; full identifiers and per-record checks are listed
in \texttt{R20-TRACE-MATRIX.csv} in S1.

Table~\ref{tab:q8} reports the pre-task process value, selected setting, and
run-tagged sample for each execution. All six samples remained outside target.
The experiment verifies the bounded harness-to-simulator execution and record
path, not process attainment or convergence. Task completion denotes a
completed software workflow, not satisfaction of the manufacturing objective.
A single probe is not a response trajectory, and the SQLite fallback does not
test the normal online DAQ backend. R20 also does not isolate
profile/binding reuse or establish cross-harness portability.

\begin{table}[!t]
\centering
\caption{Constrained model selections and simulator measurements. The control
is screw speed for cast film and the zone-3 furnace setpoint for annealing.
Each row reports one run-tagged probe after a 20-s delay from verified route
readback. All six samples are outside target; executions are sequential within
each scene.}
\label{tab:q8}
\footnotesize
\setlength{\tabcolsep}{2.5pt}
\begin{tabularx}{\columnwidth}{@{}lXcX@{}}
\toprule
Case & Pre-task value / target band & Control: start$\to$trial & Sample \\
\midrule
Film-1 & $55.82/48$--$52~\mu$m & $120\to125$~rpm & $45.69~\mu$m \\
Film-2 & $45.69/48$--$52~\mu$m & $120\to130$~rpm & $46.69~\mu$m \\
Film-3 & $46.69/48$--$52~\mu$m & $120\to125$~rpm & $45.69~\mu$m \\
Anneal-1 & $675.40/700$--$760~^{\circ}$C & $730\to745$~$^{\circ}$C & $679.60~^{\circ}$C \\
Anneal-2 & $678.30/700$--$760~^{\circ}$C & $730\to745$~$^{\circ}$C & $677.70~^{\circ}$C \\
Anneal-3 & $676.90/700$--$760~^{\circ}$C & $730\to745$~$^{\circ}$C & $678.10~^{\circ}$C \\
\bottomrule
\end{tabularx}
\end{table}

\subsection{Software Contract Checks}
A separate regression run on the local v0.8.0 working tree records 65 passing
software tests across eight suites and no failures. Denied or expired
approvals, changed recipe or limit state after approval, and invalid node
values are refused before driver actuation; unchanged approved state dispatches
through the mock driver and readback. The suites also cover tool scope, grants,
and role-specific exposure. These checks use temporary or mock state and are
separate from R20's 76 runner checks; they do not represent process trials or
test missing/mismatched DAQ run tags in the end-to-end path.

\subsection{Exploratory Protocol and Data Path (R13)}
\label{sec:crossscene}
The R13 application report describes five protocol-family read probes (OPC UA,
MQTT, Modbus TCP, Modbus RTU over TCP, and HTTP/REST) and bounded writes with
matching readbacks in two simulator scenes. Six sequential 120-s windows per
scene at 1-s sampling produced lineage-verified datasets with 658 cast-film
rows and 655 annealing rows. The rows are overlapping time-window samples, not
independent process replications. The exact runtime source was not recorded. R13 is
used only as exploratory protocol and data-lineage evidence; its archived gate
outputs and run-boundary audit are documented in the reproducibility
supplement and do not support the primary registered-harness result.

Earlier protocol, admission-ablation, controller, deployment, and scripted
recipe records in S1 use older versions or interfaces. They are not pooled
with R13 or R20, used as matched strategy comparisons, or relied on for the
current integration claims.
